
\section{3. Background}

\subsection{3.1 Softmax Attention}

For a single head with $Q, K, V \in \mathbb{R}^{L \times d_h}$:

\[
S = \frac{Q K^\top}{\sqrt{d_h}} \in \mathbb{R}^{L \times L}
\tag{3.1}
\]

\[
P = \text{softmax}(S) \in \mathbb{R}^{L \times L}
\tag{3.2}
\]

\[
O = P V \in \mathbb{R}^{L \times d_h}
\tag{3.3}
\]

\textbf{Complexity.} $\mathcal{O}(L^2)$ memory and time.

\subsection{3.2 Window Attention}

\textbf{Fixed window.} Restrict attention to positions within distance $w$:

\[
P_{ij} = 0 \quad \text{if } |i - j| > w
\tag{3.4}
\]

For a causal model, $P_{ij} = 0$ for $j > i$ or $i - j > w$.

\textbf{Complexity.} $\mathcal{O}(L w)$ memory and time.

\textbf{Implementation.} Often done with band masks or block-sparse patterns.

\subsection{3.3 Learned Span (Sukhbaatar et al., 2019)}

\textbf{Idea.} Replace the fixed window $w$ with a per-head learned span $z_h$.

\textbf{Soft ramp.} Instead of hard cutoff, use a \textbf{continuous ramp}:

\[
g_z(d) = \text{clamp}\!\left(\frac{R + z - d}{R}, 0, 1\right)
\tag{3.5}
\]

\begin{itemize}
\item For $d \leq z$: $g_z(d) = 1$ (fully attend)
\end{itemize}
\begin{itemize}
\item For $z < d < z + R$: $g_z(d) = \frac{R + z - d}{R}$ (linear decay)
\end{itemize}
\begin{itemize}
\item For $d \geq z + R$: $g_z(d) = 0$ (no attend)
\end{itemize}

\textbf{Parameterization.} We use a scalar logit:

\[
z = \sigma(\ell) \cdot S_\text{max}
\tag{3.6}
\]

\subsection{3.4 Span Regularization}

\textbf{L1 penalty.} To encourage small spans:

\[
\mathcal{L}_\text{span} = \lambda \sum_{h=1}^H z_h
\tag{3.7}
\]

At initialization, $z_h \approx S_\text{max}$ and the penalty is large. During 
training, the model shrinks $z_h$ to reduce the penalty, subject to preserving 
task accuracy. The equilibrium depends on $\lambda$: larger $\lambda$ $\to$ 
smaller spans.

\subsection{3.5 Softmax Renormalization}

After applying the ramp mask:

\[
\tilde{P}_{ij} = \frac{P_{ij} \cdot g_{z_h}(|i - j|)}{\sum_j P_{ij} \cdot g_{z_h}(|i - j|)}
\tag{3.8}
\]

This renormalization ensures rows sum to 1. We implement this via:

\[
\tilde{P} = \text{softmax}(S + \log(g_z))
\tag{3.9}
\]

where we work in log space to avoid numerical issues.

\subsection{3.6 LoRA Recovery}

LoRA (Hu et al., 2021) adds trainable low-rank adapters:

\[
W_\text{eff} = W_\text{frozen} + \frac{\alpha}{r_L} B A
\tag{3.10}
\]

with $B = 0$ at initialization. ASA-Span uses LoRA to adapt the frozen 
projections to the new span structure.

\subsection{3.7 Relative Position Bias}

T5 (Raffel et al., 2020) adds a learned bias based on bucketed relative 
position:

\[
S_{ij} \leftarrow S_{ij} + b_{\text{bucket}(i - j)}
\tag{3.11}
\]

ASA-Span preserves this bias, applying the ramp \textbf{after} adding the bias.

\subsection{3.8 Notation Summary}
